\documentclass[10pt,a4paper]{amsart}
\usepackage{amsmath,amssymb,amsthm,graphicx}
% Retained minimal preamble for the Stacks Project chapter more-morphisms.tex (More on Morphisms).
% Verbatim excerpt of the chapter preamble: Xy-pic setup (native preamble lines 14-18) and the native
% theorem declarations with the mathematical macros (native preamble lines 154-217). Copied mechanically;
% nothing added, removed or redefined. The wrapper class and the amsmath/amssymb/amsthm/graphicx base
% packages are supplied by the assembler, so the class/externaldocument part of the native preamble is omitted.
% For commutative diagrams we use Xy-pic
\usepackage[all]{xy}

% We use 2cell for 2-commutative diagrams.
\xyoption{2cell}
% Theorem environments.
%
\theoremstyle{plain}
\newtheorem{theorem}[subsection]{Theorem}
\newtheorem{proposition}[subsection]{Proposition}
\newtheorem{lemma}[subsection]{Lemma}

\theoremstyle{definition}
\newtheorem{definition}[subsection]{Definition}
\newtheorem{example}[subsection]{Example}
\newtheorem{exercise}[subsection]{Exercise}
\newtheorem{situation}[subsection]{Situation}

\theoremstyle{remark}
\newtheorem{remark}[subsection]{Remark}
\newtheorem{remarks}[subsection]{Remarks}

\numberwithin{equation}{subsection}

% Macros
%
\def\lim{\mathop{\mathrm{lim}}\nolimits}
\def\colim{\mathop{\mathrm{colim}}\nolimits}
\def\Spec{\mathop{\mathrm{Spec}}}
\def\Hom{\mathop{\mathrm{Hom}}\nolimits}
\def\Ext{\mathop{\mathrm{Ext}}\nolimits}
\def\SheafHom{\mathop{\mathcal{H}\!\mathit{om}}\nolimits}
\def\SheafExt{\mathop{\mathcal{E}\!\mathit{xt}}\nolimits}
\def\Sch{\mathit{Sch}}
\def\Mor{\mathop{\mathrm{Mor}}\nolimits}
\def\Ob{\mathop{\mathrm{Ob}}\nolimits}
\def\Sh{\mathop{\mathit{Sh}}\nolimits}
\def\NL{\mathop{N\!L}\nolimits}
\def\CH{\mathop{\mathrm{CH}}\nolimits}
\def\proetale{{pro\text{-}\acute{e}tale}}
\def\etale{{\acute{e}tale}}
\def\QCoh{\mathit{QCoh}}
\def\Ker{\mathop{\mathrm{Ker}}}
\def\Im{\mathop{\mathrm{Im}}}
\def\Coker{\mathop{\mathrm{Coker}}}
\def\Coim{\mathop{\mathrm{Coim}}}

% Boxtimes
%
\DeclareMathSymbol{\boxtimes}{\mathbin}{AMSa}{"02}

%
% Macros for moduli stacks/spaces
%
\def\QCohstack{\mathcal{QC}\!\mathit{oh}}
\def\Cohstack{\mathcal{C}\!\mathit{oh}}
\def\Spacesstack{\mathcal{S}\!\mathit{paces}}
\def\Quotfunctor{\mathrm{Quot}}
\def\Hilbfunctor{\mathrm{Hilb}}
\def\Curvesstack{\mathcal{C}\!\mathit{urves}}
\def\Polarizedstack{\mathcal{P}\!\mathit{olarized}}
\def\Complexesstack{\mathcal{C}\!\mathit{omplexes}}
% \Pic is the operator that assigns to X its picard group, usage \Pic(X)
% \Picardstack_{X/B} denotes the Picard stack of X over B
% \Picardfunctor_{X/B} denotes the Picard functor of X over B
\def\Pic{\mathop{\mathrm{Pic}}\nolimits}
\def\Picardstack{\mathcal{P}\!\mathit{ic}}
\def\Picardfunctor{\mathrm{Pic}}
\def\Deformationcategory{\mathcal{D}\!\mathit{ef}}
\newenvironment{reference}{}{}
\newenvironment{slogan}{}{}
\newenvironment{history}{}{}
\usepackage{cleveref}
\crefname{originallemma}{original lemma}{original lemmas}
\Crefname{originallemma}{Original lemma}{Original lemmas}
\makeatletter
\newlabel{descent-lemma-CM-local-syntomic}{{036B}{source}{}{sourcepointer-1}{}}
\newlabel{descent-lemma-CM-local-syntomic@cref}{{[originallemma][0][]036B}{[1][0][]source}{}{}{}}
\makeatother
\crefname{originallemma}{original lemma}{original lemmas}
\Crefname{originallemma}{Original lemma}{Original lemmas}
\makeatletter
\newlabel{descent-lemma-descending-property-flat}{{02L2}{source}{}{sourcepointer-2}{}}
\newlabel{descent-lemma-descending-property-flat@cref}{{[originallemma][0][]02L2}{[1][0][]source}{}{}{}}
\makeatother
\crefname{originallemma}{original lemma}{original lemmas}
\Crefname{originallemma}{Original lemma}{Original lemmas}
\makeatletter
\newlabel{descent-lemma-flat-fpqc-local-source}{{036K}{source}{}{sourcepointer-3}{}}
\newlabel{descent-lemma-flat-fpqc-local-source@cref}{{[originallemma][0][]036K}{[1][0][]source}{}{}{}}
\makeatother
\crefname{originallemma}{original lemma}{original lemmas}
\Crefname{originallemma}{Original lemma}{Original lemmas}
\makeatletter
\newlabel{varieties-lemma-CM-base-change}{{045P}{source}{}{sourcepointer-4}{}}
\newlabel{varieties-lemma-CM-base-change@cref}{{[originallemma][0][]045P}{[1][0][]source}{}{}{}}
\makeatother
\crefname{originallemma}{original lemma}{original lemmas}
\Crefname{originallemma}{Original lemma}{Original lemmas}
\makeatletter
\newlabel{varieties-lemma-locally-Noetherian-base-change}{{038R}{source}{}{sourcepointer-5}{}}
\newlabel{varieties-lemma-locally-Noetherian-base-change@cref}{{[originallemma][0][]038R}{[1][0][]source}{}{}{}}
\makeatother
\title{Cohen-Macaulayness with locally Noetherian fibres is fppf local on the target and syntomic local on the source}
\author{The Stacks Project Authors}
\date{stacks-project:more-morphisms.tex:local-snapshot-20261011; GNU FDL 1.2 or later}
\begin{document}
\maketitle
\noindent\textit{Editorial note.} Statement span is lines 6101-6107 of sources/stacks-more-morphisms/pinned-original-source.tex; proof span is lines 6109-6148, the complete author proof (38 body lines).\par\medskip
\noindent\textit{Editorial note.} Navigation stubs: 5 cited-but-not-included labels are supplied as non-load-bearing pointers with their pinned tags where the snapshot tag index has them.\par\medskip
\noindent\textit{Editorial note.} extra\_packages is empty on purpose: the retained preamble already loads xy with options and a second declaration clashes in the assembler wrapper.\par\medskip
\noindent\textit{Editorial note.} Text is retained verbatim, including the author's own typos and placeholders.\par\medskip
\noindent\textit{Allowed standard prerequisites (editorial).}\begin{itemize}
\item Explicit prerequisite boundary. Cohen-Macaulay morphisms, fppf descent, and syntomic base change; the reductions are retained.
\end{itemize}\medskip
\ifdefined\hypertarget\hypertarget{sourcepointer-1}{}\fi
\noindent\textit{Source reference.} Original lemma 036B: sources/stacks-more-morphisms/pinned-original-source.tex, source lines 6148--6148. Cited by the retained statement or proof but not part of the collected spans; supplied only as a navigation pointer. The argument does not depend on its content.\par\medskip
\ifdefined\hypertarget\hypertarget{sourcepointer-2}{}\fi
\noindent\textit{Source reference.} Original lemma 02L2: sources/stacks-more-morphisms/pinned-original-source.tex, source lines 6148--6148. Cited by the retained statement or proof but not part of the collected spans; supplied only as a navigation pointer. The argument does not depend on its content.\par\medskip
\ifdefined\hypertarget\hypertarget{sourcepointer-3}{}\fi
\noindent\textit{Source reference.} Original lemma 036K: sources/stacks-more-morphisms/pinned-original-source.tex, source lines 6148--6148. Cited by the retained statement or proof but not part of the collected spans; supplied only as a navigation pointer. The argument does not depend on its content.\par\medskip
\ifdefined\hypertarget\hypertarget{sourcepointer-4}{}\fi
\noindent\textit{Source reference.} Original lemma 045P: sources/stacks-more-morphisms/pinned-original-source.tex, source lines 6148--6148. Cited by the retained statement or proof but not part of the collected spans; supplied only as a navigation pointer. The argument does not depend on its content.\par\medskip
\ifdefined\hypertarget\hypertarget{sourcepointer-5}{}\fi
\noindent\textit{Source reference.} Original lemma 038R: sources/stacks-more-morphisms/pinned-original-source.tex, source lines 6148--6148. Cited by the retained statement or proof but not part of the collected spans; supplied only as a navigation pointer. The argument does not depend on its content.\par\medskip
\begin{lemma}
\label{lemma-CM-local-source-and-target}
The property
$\mathcal{P}(f)=$``the fibres of $f$ are locally Noetherian and $f$ is
Cohen-Macaulay'' is local in the fppf topology on the target and
local in the syntomic topology on the source.
\end{lemma}

\begin{proof}
We have
$\mathcal{P}(f) =
\mathcal{P}_1(f) \wedge \mathcal{P}_2(f)$
where
$\mathcal{P}_1(f)=$``$f$ is flat'', and
$\mathcal{P}_2(f)=$``the fibres of $f$ are locally Noetherian
and Cohen-Macaulay''.
We know that $\mathcal{P}_1$ is
local in the fppf topology on the source and the target, see
Descent, Lemmas \ref{descent-lemma-descending-property-flat} and
\ref{descent-lemma-flat-fpqc-local-source}. Thus we have to deal
with $\mathcal{P}_2$.

\medskip\noindent
Let $f : X \to Y$ be a morphism of schemes.
Let $\{\varphi_i : Y_i \to Y\}_{i \in I}$ be an fppf covering of $Y$.
Denote $f_i : X_i \to Y_i$ the base change of $f$ by $\varphi_i$.
Let $i \in I$ and let $y_i \in Y_i$ be a point.
Set $y = \varphi_i(y_i)$. Note that
$$
X_{i, y_i} = \Spec(\kappa(y_i)) \times_{\Spec(\kappa(y))} X_y.
$$
and that $\kappa(y_i)/\kappa(y)$ is a finitely generated field
extension. Hence if $X_y$ is locally Noetherian, then
$X_{i, y_i}$ is locally Noetherian, see
Varieties, Lemma \ref{varieties-lemma-locally-Noetherian-base-change}.
And if in addition $X_y$ is Cohen-Macaulay,
then $X_{i, y_i}$ is Cohen-Macaulay, see
Varieties, Lemma \ref{varieties-lemma-CM-base-change}.
Thus $\mathcal{P}_2$ is fppf local on the target.

\medskip\noindent
Let $\{X_i \to X\}$ be a syntomic covering of $X$.
Let $y \in Y$. In this case $\{X_{i, y} \to X_y\}$ is a
syntomic covering of the fibre. Hence the locality of $\mathcal{P}_2$
for the syntomic topology on the source follows from
Descent, Lemma \ref{descent-lemma-CM-local-syntomic}.
Combining the above the lemma follows.
\end{proof}

\end{document}
